\section{A criterion for finitely many row types}\label{sec:row-criterion}

This section contains the common core of the proofs of
Theorems~\ref{thm:upper-tail} and~\ref{thm:fixed-densities}.  It is stated
without reference to either variational problem.  The input is a functional on
rows whose global minimum is attained by the rows of a graphon $W$; the output
is that $W$ is multipodal with ordered blocks.

Throughout this section, we fix a graph $H$ with $\Delta(H)\ge2$, a graphon
$W$, and numbers $q\in(0,1)$, $\alpha\in\RR$, $\mu>0$, and $\omega>0$ such that
\begin{equation}\label{eq:positive-base}
  W\ge\omega\qquad\text{almost everywhere}.
\end{equation}
For measurable $u:[0,1]\to[0,1]$, define the \emph{row functional}
\begin{equation}\label{eq:row-functional}
  \mathcal F_W(u):=2\int_0^1J_q\bigl(u(y)\bigr)\dd y-2\alpha\int_0^1u(y)\dd y
  -\mu\,\mathcal R_W(u).
\end{equation}
Assume that there is $c_W\in\RR$ such that
\begin{equation}\label{eq:row-lower-bound}
  \mathcal F_W(u)\ge c_W\qquad\text{for every measurable }u:[0,1]\to[0,1],
\end{equation}
and let
\[
  \mathcal M_W:=\bigl\{u\in L^\infty([0,1]):0\le u\le1,\ \mathcal F_W(u)=c_W\bigr\}
\]
be the set of global minimizers of $\mathcal F_W$ attaining this bound.

\begin{theorem}[Finitely many row types]\label{thm:row-criterion}
Under the assumptions above, the set $\mathcal M_W$ is finite, and any two of
its elements $u_1,u_2$ satisfy $u_1\le u_2$ or $u_2\le u_1$ almost everywhere.
If, in addition,
\begin{equation}\label{eq:rows-attain}
  W_x\in\mathcal M_W\qquad\text{for almost every }x\in[0,1],
\end{equation}
then $W$ is multipodal and its blocks are ordered.
\end{theorem}

The proof has three steps.  First, the singularities of $J_q'$ at $0$ and $1$
confine all minimizers to a common interval $[\delta,1-\delta]$, where they
satisfy an Euler--Lagrange equation.  Second, strict supermodularity of
$\mathcal R_W$ shows that two minimizers cannot cross, so $\mathcal M_W$ is
totally ordered.  Third, compactness and an analysis of the second variation
of $\mathcal F_W$ show that a totally ordered set of minimizers cannot contain
a continuum, and is therefore finite.  The assumption~\eqref{eq:positive-base}
is used only in the strict inequalities of the second and third steps.

\subsection{Interiority and the Euler--Lagrange equation}

\begin{lemma}[Interiority]\label{lem:row-interior}
There is $\delta\in(0,1/2)$ such that every $u\in\mathcal M_W$ satisfies
$\delta\le u\le1-\delta$.  Moreover, every $u\in\mathcal M_W$ satisfies
\begin{equation}\label{eq:row-euler}
  2J_q'(u)-2\alpha=\mu\nabla\mathcal R_W(u),
  \qquad\text{equivalently}\qquad
  u=\sigma\Bigl(\lo(q)+\alpha+\frac\mu2\nabla\mathcal R_W(u)\Bigr).
\end{equation}
\end{lemma}

\begin{proof}
Choose $\delta\in(0,1/2)$ so small that $2J_q'(z)-2\alpha+2m\mu<0$ for
$0<z<\delta$ and $2J_q'(z)-2\alpha-2m\mu>0$ for $1-\delta<z<1$.  Let
$u\in\mathcal M_W$ and suppose that $A:=\{u<\delta\}$ has positive measure.
Put $\tilde u:=\max(u,\delta)$.  By~\eqref{eq:rooted-lipschitz},
\[
  \mathcal F_W(\tilde u)-\mathcal F_W(u)
  \le\int_A\Bigl(2\bigl(J_q(\delta)-J_q(u)\bigr)-2\alpha(\delta-u)
  +2m\mu(\delta-u)\Bigr)
  =\int_A\int_u^\delta\bigl(2J_q'(z)-2\alpha+2m\mu\bigr)\dd z<0,
\]
contradicting~\eqref{eq:row-lower-bound}.  The case $\{u>1-\delta\}$ is
analogous, with $\tilde u:=\min(u,1-\delta)$.  Since every $u\in\mathcal M_W$
takes values in $[\delta,1-\delta]$, the functions $u+th$ with bounded $h$ and
small $|t|$ are admissible, so the first variation of $\mathcal F_W$ at $u$
vanishes in every bounded direction.  This gives~\eqref{eq:row-euler}.
\end{proof}

\subsection{Minimizers cannot cross}

\begin{lemma}[Strict supermodularity]\label{lem:strict-supermodular}
Let $u_1,u_2\in\mathcal M_W$.  Then
\[
  \mathcal R_W(u_1\wedge u_2)+\mathcal R_W(u_1\vee u_2)
  \ge\mathcal R_W(u_1)+\mathcal R_W(u_2),
\]
and the inequality is strict if $u_1$ and $u_2$ cross, that is, if both
$\{u_1>u_2\}$ and $\{u_2>u_1\}$ have positive measure.
\end{lemma}

\begin{proof}
Put $g:=u_1\wedge u_2$, $h_1:=(u_1-u_2)_+$, and $h_2:=(u_2-u_1)_+$, so that
$u_1=g+h_1$, $u_2=g+h_2$, and $u_1\vee u_2=g+h_1+h_2$.  Expanding the products
in~\eqref{eq:rooted-density}, the difference of the two sides is a sum of
integrals of nonnegative kernels against products containing at least one
factor $h_1$ and at least one factor $h_2$.  This proves the weak inequality.
For strictness, choose a vertex $v_0$ with degree $d_0\ge2$.  Among the terms
of $\mathcal R_{v_0,W}$, keep only those with exactly one factor $h_1$, exactly
one factor $h_2$, and $d_0-2$ factors $g$.  Each edge factor not incident to
$v_0$ is at least $\omega$ by~\eqref{eq:positive-base}, and $g\ge\delta$ by
Lemma~\ref{lem:row-interior}, so
\[
  \mathcal R_W(u_1\wedge u_2)+\mathcal R_W(u_1\vee u_2)
  -\mathcal R_W(u_1)-\mathcal R_W(u_2)
  \ge d_0(d_0-1)\,\omega^{m-d_0}\delta^{d_0-2}
  \int_0^1h_1\int_0^1h_2,
\]
which is positive when $u_1$ and $u_2$ cross.
\end{proof}

\begin{prop}[Total order]\label{prop:row-chain}
For all $u_1,u_2\in\mathcal M_W$, either $u_1\le u_2$ or $u_2\le u_1$ almost
everywhere.
\end{prop}

\begin{proof}
Pointwise, the two numbers $(u_1\wedge u_2)(y)$ and $(u_1\vee u_2)(y)$ are
$u_1(y)$ and $u_2(y)$ in some order.  Hence the sum of the first two terms
of~\eqref{eq:row-functional} over the pair $u_1\wedge u_2$, $u_1\vee u_2$ equals
the corresponding sum over the pair $u_1$, $u_2$.  If $u_1$ and $u_2$ crossed,
Lemma~\ref{lem:strict-supermodular} and $\mu>0$ would give
\[
  \mathcal F_W(u_1\wedge u_2)+\mathcal F_W(u_1\vee u_2)
  <\mathcal F_W(u_1)+\mathcal F_W(u_2)=2c_W,
\]
contradicting~\eqref{eq:row-lower-bound} applied to $u_1\wedge u_2$ and
$u_1\vee u_2$.
\end{proof}

\subsection{Compactness and the second variation}

\begin{lemma}[Compactness]\label{lem:row-compact}
The set $\mathcal M_W$ is compact in $L^\infty([0,1])$.
\end{lemma}

\begin{proof}
Each term~\eqref{eq:rooted-gradient} of $\nabla\mathcal R_W(u)$ has the form
$y\mapsto\int K(y,z_1,\dots,z_k)\prod_{i=1}^ku(z_i)\dd z$ with a measurable
kernel $0\le K\le1$.  Such a kernel defines a Hilbert--Schmidt, hence compact,
operator from $L^2([0,1]^k)$ to $L^2([0,1])$, and the functions
$u^{\otimes k}$ with $0\le u\le1$ form a bounded subset of $L^2([0,1]^k)$.
Hence $\{\nabla\mathcal R_W(u):0\le u\le1\}$ is relatively compact in $L^2$.
Moreover, expanding the difference of products telescopically gives
\begin{equation}\label{eq:gradient-lipschitz}
  \|\nabla\mathcal R_W(u)-\nabla\mathcal R_W(u')\|_\infty\le C_H\|u-u'\|_1.
\end{equation}
Now let $(u_j)\subseteq\mathcal M_W$.  After passing to a subsequence,
$\nabla\mathcal R_W(u_j)\to g$ in $L^2$.  Since $\sigma$ is Lipschitz, the
fixed-point form of~\eqref{eq:row-euler} gives
$u_j\to u:=\sigma(\lo(q)+\alpha+\frac\mu2g)$ in $L^2$.
By~\eqref{eq:gradient-lipschitz}, $\nabla\mathcal R_W(u_j)\to\nabla\mathcal R_W(u)$
uniformly, so $g=\nabla\mathcal R_W(u)$, and~\eqref{eq:row-euler} now gives
$u_j\to u$ uniformly.  Since $\mathcal F_W$ is continuous under uniform
convergence, $u\in\mathcal M_W$.
\end{proof}

For $u\in\mathcal M_W$, consider the second variation
\[
  Q_u(h):=2\int_0^1J_q''(u)\,h^2-\mu\,D^2\mathcal R_W(u)[h,h],
  \qquad h\in L^2([0,1]).
\]
The bilinear form $D^2\mathcal R_W(u)$ is given by a bounded symmetric kernel,
so it is represented by a compact self-adjoint operator $B_u$ on
$L^2([0,1])$.  By Lemma~\ref{lem:row-interior}, multiplication by
$2J_q''(u)=2/(u(1-u))$ is a bounded operator $A_u$ with $A_u\ge8$ and bounded
inverse.  Thus $Q_u(h)=\langle h,\mathcal L_uh\rangle$ with
$\mathcal L_u:=A_u-\mu B_u$.

\begin{lemma}[Coercivity on mean-zero functions]\label{lem:coercive}
For every $u\in\mathcal M_W$, there is $c_u>0$ such that
$Q_u(h)\ge c_u\|h\|_2^2$ for every $h\in L^2([0,1])$ with $\int_0^1h=0$.
\end{lemma}

\begin{proof}
Since $u$ is a global minimizer and takes values in $[\delta,1-\delta]$, the
second variation is nonnegative: $Q_u(h)\ge0$ for bounded $h$, and by density
for all $h\in L^2$.  Let $h\ne0$ have mean zero, and write $h=h_+-h_-$ with
$h_\pm\ge0$ of disjoint supports; both are nonzero.  Choosing a vertex $v_0$ of
degree $d_0\ge2$ as in the proof of Lemma~\ref{lem:strict-supermodular}, we get
$D^2\mathcal R_W(u)[h_+,h_-]\ge d_0(d_0-1)\omega^{m-d_0}\delta^{d_0-2}
\|h_+\|_1\|h_-\|_1>0$.  Since $h_+$ and $h_-$ have disjoint supports,
\[
  Q_u(h)=Q_u(h_+)+Q_u(h_-)+2\mu\,D^2\mathcal R_W(u)[h_+,h_-]>0.
\]
Suppose that the quantitative bound fails.  Then there are $h_j$ with mean
zero, $\|h_j\|_2=1$, and $Q_u(h_j)\to0$.  Passing to a subsequence,
$h_j\rightharpoonup h$ weakly in $L^2$, and $h$ has mean zero.  Compactness of
$B_u$ gives $\langle h_j,B_uh_j\rangle\to\langle h,B_uh\rangle$.  If $h=0$, then
$Q_u(h_j)\ge8-\mu\langle h_j,B_uh_j\rangle\to8$, a contradiction.  If $h\ne0$,
then weak lower semicontinuity of $h\mapsto\langle h,A_uh\rangle$ gives
$Q_u(h)\le\liminf Q_u(h_j)=0$, contradicting strict positivity.
\end{proof}

\begin{lemma}[Bordered second variation]\label{lem:bordered}
For every $u\in\mathcal M_W$, the operator
$(h,\lambda)\mapsto\bigl(\mathcal L_uh-\lambda\mathbf 1,\int_0^1h\bigr)$ is a
bounded isomorphism of $L^2([0,1])\times\RR$ onto itself and of
$L^\infty([0,1])\times\RR$ onto itself.
\end{lemma}

\begin{proof}
Let $L^2_0$ be the space of mean-zero functions and $P_0$ the orthogonal
projection onto it.  By Lemma~\ref{lem:coercive} and the Lax--Milgram theorem,
$P_0\mathcal L_u$ restricted to $L^2_0$ is invertible.  Given
$(g,\beta)\in L^2\times\RR$, write $h=\beta\mathbf1+k$ with $k\in L^2_0$, solve
$P_0\mathcal L_uk=P_0(g-\mathcal L_u(\beta\mathbf 1))$, and put
$\lambda:=\int_0^1(\mathcal L_uh-g)$; then $\mathcal L_uh-\lambda\mathbf1=g$.  If
$(h,\lambda)$ is in the kernel, then $h\in L^2_0$ and
$Q_u(h)=\lambda\int h=0$, so $h=0$ and $\lambda=0$.  This proves the claim on
$L^2\times\RR$.  If $g\in L^\infty$, then
$h=A_u^{-1}(g+\lambda\mathbf1+\mu B_uh)$, and $B_u$ maps $L^2$ boundedly into
$L^\infty$ because its kernel is bounded; hence $h\in L^\infty$ with a
corresponding norm bound.
\end{proof}

\subsection{Finiteness}

\begin{lemma}[Finiteness]\label{lem:row-finite}
The set $\mathcal M_W$ is finite.
\end{lemma}

\begin{proof}
We may assume $\mathcal M_W\ne\emptyset$.  By Proposition~\ref{prop:row-chain},
distinct elements of $\mathcal M_W$ differ on a set of positive measure with a
fixed sign, so they have distinct means.  Hence the map
$\mathfrak m(u):=\int_0^1u$ is injective on $\mathcal M_W$, and by
Lemma~\ref{lem:row-compact} it is a homeomorphism from $\mathcal M_W$ onto the
compact set $\mathcal A:=\mathfrak m(\mathcal M_W)\subseteq\RR$.

Fix $u_*\in\mathcal M_W$ and put $a_*:=\mathfrak m(u_*)$.  Consider the map
\[
  \Xi(a;u,\lambda):=\Bigl(2J_q'(u)-2\alpha-\mu\nabla\mathcal R_W(u)-\lambda\mathbf1,\
  \int_0^1u-a\Bigr),
\]
defined for $a\in\RR$, $\lambda\in\RR$, and $u$ in the open $L^\infty$-ball of
radius $\delta/2$ around $u_*$.  It is real-analytic with values in
$L^\infty\times\RR$: the map $u\mapsto J_q'(u)$ is analytic on this ball, and
$u\mapsto\nabla\mathcal R_W(u)$ is polynomial.  We have $\Xi(a_*;u_*,0)=0$ by
\eqref{eq:row-euler}, and the derivative of $\Xi$ in $(u,\lambda)$ is the
isomorphism of Lemma~\ref{lem:bordered}.  By Lemma~\ref{lem:analytic-ift},
there are an open interval $I_*\ni a_*$, a neighborhood of $(u_*,0)$, and
analytic maps $a\mapsto u_*(a)$ and $a\mapsto\lambda_*(a)$ on $I_*$ such that,
in this neighborhood, $\Xi(a;u,\lambda)=0$ if and only if
$(u,\lambda)=(u_*(a),\lambda_*(a))$.

Suppose that $a_*$ is an accumulation point of $\mathcal A$, and choose distinct
$a_j\in\mathcal A$ with $a_j\to a_*$ and $u_j\in\mathcal M_W$ with
$\mathfrak m(u_j)=a_j$.  Then $u_j\to u_*$ in $L^\infty$, and each $u_j$ satisfies
$\Xi(a_j;u_j,0)=0$ by~\eqref{eq:row-euler}.  By local uniqueness,
$\lambda_*(a_j)=0$ for all large $j$.  The analytic function $\lambda_*$ has
zeros accumulating at $a_*$, so $\lambda_*\equiv0$ near $a_*$.  Along the
branch, differentiating $\int u_*(a)=a$ gives
\[
  \frac{\dd}{\dd a}\mathcal F_W\bigl(u_*(a)\bigr)
  =\int_0^1\bigl(2J_q'(u_*(a))-2\alpha-\mu\nabla\mathcal R_W(u_*(a))\bigr)u_*'(a)
  =\lambda_*(a)=0.
\]
Hence $\mathcal F_W(u_*(a))=c_W$ near $a_*$, so $u_*(a)\in\mathcal M_W$ and a
whole open interval around $a_*$ lies in $\mathcal A$.

We have shown that every accumulation point of $\mathcal A$ is an interior
point of $\mathcal A$.  If $\mathcal A$ were infinite, its set $\mathcal A'$ of
accumulation points would be nonempty and compact.  By what we have shown,
every point of $\mathcal A'$ has a neighborhood contained in $\mathcal A$, all
of whose points are accumulation points; hence $\mathcal A'$ is open.  A
nonempty subset of $\RR$ cannot be both open and compact, so $\mathcal A$, and
therefore $\mathcal M_W$, is finite.
\end{proof}

\begin{proof}[Proof of Theorem~\ref{thm:row-criterion}]
Proposition~\ref{prop:row-chain} and Lemma~\ref{lem:row-finite} show that
$\mathcal M_W=\{g_1,\dots,g_M\}$ is a finite totally ordered set; label its
elements so that $g_1\le\dots\le g_M$.  Assume~\eqref{eq:rows-attain}, and let
$A_i:=\{x\in[0,1]:W_x=g_i\text{ almost everywhere}\}$.  These sets are
measurable, pairwise disjoint, and cover $[0,1]$ up to a null set; we discard
the null sets among them and add the remaining null set to one of the others.
For almost every $(x,y)\in A_i\times A_j$, the row identities and the symmetry
of $W$ give $g_i(y)=W(x,y)=W(y,x)=g_j(x)$.  A function of $y$ that agrees almost
everywhere on $A_i\times A_j$ with a function of $x$ is almost everywhere
constant, so $W=b_{ij}$ almost everywhere on $A_i\times A_j$ for some constant
$b_{ij}=b_{ji}$.  Thus $W$ is multipodal.  Finally, $g_i=b_{ij}$ on $A_j$, so
$g_i\le g_{i'}$ for $i\le i'$ is exactly~\eqref{eq:ordered-blocks}.
\end{proof}

\begin{rmk}\label{rmk:criterion-input}
Theorem~\ref{thm:row-criterion} does not use that $W$ solves any graphon
variational problem, nor an Euler--Lagrange equation for $W$ itself.  All
information about the ambient problem enters through the global lower
bound~\eqref{eq:row-lower-bound}, the sign $\mu>0$, the lower
bound~\eqref{eq:positive-base}, and the property~\eqref{eq:rows-attain}.  In
Sections~\ref{sec:upper-tail} and~\ref{sec:fixed-densities}, these hypotheses
are supplied by inserting a small block of vertices into an optimizer.
\end{rmk}
